\documentclass[12pt]{article}

% --- Packages ---
\usepackage[utf8]{inputenc}
\usepackage{geometry}
\usepackage{amsmath, amssymb}
\usepackage{hyperref}
\usepackage{graphicx} % Required for \includegraphics
\usepackage{fancyhdr}
\usepackage{enumitem}
\usepackage{autobreak}
\usepackage{amsmath}
\usepackage{breqn}

\usepackage{amsmath,amssymb}
\usepackage{parskip}       % Correct paragraph spacing
\usepackage{xcolor}        % Required for \textcolor command
\usepackage{enumitem}
\usepackage{float}
\usepackage{listings}
\usepackage{xcolor}
\usepackage{tikz}
\usepackage[utf8]{inputenc}
\usepackage[T1]{fontenc}
\usepackage{url}
\usepackage{amsmath,amssymb,amsthm}
\usepackage{booktabs}
\usepackage{pgfplotstable}
\usepackage{siunitx}
\usepackage{geometry}

\setlist{nosep}           % Reduce vertical spacing in lists
\usepackage{amsmath,amssymb}
\usepackage{geometry}
\usepackage{listings}
\usepackage{booktabs}
\usepackage{graphicx}
\usepackage{listings}
\usepackage{pgfplots}
\pgfplotsset{compat=1.18} % sau ultima versiune instalată la tine
\usepackage{filecontents} % dacă 
% --- Page Setup ---
\geometry{a4paper, margin=1in, top=1.2in, headheight=35pt} % Expanded top margin & headheight for the ornament

\lstset{
    language=python,
    basicstyle=\ttfamily\small,
    backgroundcolor=\color{blue!5},
    keywordstyle=\color{purple!80!blue},
    commentstyle=\color{green!50!black},
    stringstyle=\color{red!70!black},
    numberstyle=\tiny\color{gray},
    numbers=left,
    numbersep=15pt,
    frame=shadowbox,
    rulesepcolor=\color{blue!20},
    frameround=tttt,
    breaklines=true,
    captionpos=b,
    showstringspaces=false,
    tabsize=4,
    xleftmargin=20pt,
    xrightmargin=5pt
}

% --- Header & Footer Configuration ---
\pagestyle{fancy}
\fancyhf{} % Clear default headers/footers

% Center the ornament in the header (antet)
\chead{\includegraphics[width=4cm]{ornament_6.pdf}}

% Center the page number in the footer
\cfoot{\thepage}

% Optional: Remove the default rule line below the header
\renewcommand{\headrulewidth}{0pt}

\begin{document}

\begin{titlepage}
    \thispagestyle{fancy} % Force header/footer onto the title page
    \centering
    
    \vspace*{4cm}

    % Main Title
    {\Huge \bfseries Homework 1}

    \vspace{4cm}

    % Name and Date
    {\Large \textbf{Name:} Gatman Vitalie} \\[0.5em]
    {\Large \textbf{Date:} \today} \\[0.5em]
    {\Large \textbf{Group:} FAF-263}

    \vfill
\end{titlepage}

\newpage

% Page 2 content...
\section*{Problem 1.3}

We define the power set $2^A$ as the set containing all possible distinct
subsets of $A$, including the empty set $\varnothing$ and $A$ itself.
For the given examples, we have:

\begin{enumerate}[label=\alph*)]

    \item For $A = \{2,\,5\}$, we have
    \[
    2^A = \left\{
    \varnothing,\,
    \{2\},\,
    \{5\},\,
    \{2,\,5\}
    \right\}.
    \]

    \item For $A = \{f,\,a,\,f\}$, we have
    \[
    2^A = \left\{
    \varnothing,\,
    \{f\},\,
    \{a\},\,
    \{f,\,a\}
    \right\}.
    \]
    This is because sets cannot contain duplicate elements. Therefore,
    \[
    A = \{f,\,a,\,f\} = \{f,\,a\},
    \]
    so $2^A = 2^{\{f,\,a\}}$.

    \item For $A = \{\varnothing,\,\{z\}\}$, we have
    \[
    2^A =
    \left\{
    \varnothing,\,
    \{\varnothing\},\,
    \{\{z\}\},\,
    \{\varnothing,\,\{z\}\}
    \right\}.
    \]
    This is because $A$ has two distinct elements, namely $\varnothing$
    and $\{z\}$. Therefore, its power set contains all four possible
    subsets of these two elements.

\end{enumerate}

To find the power set of the power set of $A$, $2^{2^A}$, we first compute
$2^A$. We then treat $2^A$ as an ordinary set and compute its power set.
For instance, in example (a), we have $A = \{2,\,5\}$. If we define
$B = 2^A$, then

\[
B = \left\{
\varnothing,\,
\{2\},\,
\{5\},\,
\{2,\,5\}
\right\}.
\]

Therefore,

\[
2^{2^A} = 2^B,
\]

which is simply the power set of $B$.
\vspace{1cm}
\subsection*{My code for iterating power sets}
\begin{lstlisting}[language=python]
from itertools import combinations

def latex_element(element):
    """Convert a single element to LaTeX."""

    if element == EMPTY_SET:
        return r"\varnothing"

    if isinstance(element, str):
        if element.strip() in {"\\varnothing"}:
            return r"\varnothing"

        replacements = {
            "\\": r"\textbackslash{}",
            "&": r"\&",
            "%": r"\%",
            "$": r"\$",
            "#": r"\#",
            "_": r"\_",
            "{": r"\{",
            "}": r"\}",
        }

        result = element
        for old, new in replacements.items():
            result = result.replace(old, new)

        return r"\text{" + result + "}"

    return str(element)


def latex_set(s):
    """Recursively convert a Python set/frozenset into LaTeX."""

    if len(s) == 0:
        return r"\varnothing"

    elements = sorted(s, key=str)

    formatted = []

    for element in elements:

        if isinstance(element, (set, frozenset)):
            formatted.append(latex_set(element))

        elif element == EMPTY_SET:
            formatted.append(r"\varnothing")

        else:
            formatted.append(latex_element(element))

    return r"\{" + r",\, ".join(formatted) + r"\}"


# ---------- Power set ----------

def power_set(s):
    """Return the power set of s as a set of frozensets."""

    elements = list(s)

    return {
        frozenset(combination)
        for r in range(len(elements) + 1)
        for combination in combinations(elements, r)
    }


# ---------- Input ----------

EMPTY_SET = "__EMPTY_SET__"


def parse_element(x):

    x = x.strip()

    if x in {"\\varnothing"}:
        return EMPTY_SET

    try:
        return int(x)
    except ValueError:
        pass

    try:
        return float(x)
    except ValueError:
        pass

    return x


def main():

    print("Power Set Generator")
    print("--------------------")

    raw = input(
        "Enter the elements separated by commas "
        "(example: 2, 5, hello, world):\n> "
    )

    elements = [parse_element(x) for x in raw.split(",")]

    # Remove dopubles
    original_set = frozenset(elements)

    iterations = int(
        input(
            "\nNumber of iterations "
            "(1 = P(A), 2 = P(P(A)), etc.):\n> "
        )
    )

    filename = "powerset.tex"

    if not filename:
        filename = "power_set.tex"

    if not filename.endswith(".tex"):
        filename += ".tex"

    current = original_set

    for i in range(iterations):
        current = power_set(current)

        print(
            f"Iteration {i + 1}: "
            f"{len(current)} elements"
        )

    output = latex_set(current)

    with open(filename, "w", encoding="utf-8") as file:
        file.write(output)

    print(f"\nDone!")
    print(f"Result written to: {filename}")

    if len(output) < 5000:
        print("\nLaTeX output:")
        print(output)
    else:
        print(
            f"\nThe LaTeX output is {len(output):,} characters long."
        )


if __name__ == "__main__":
    main()
\end{lstlisting}

\newpage


\section*{Problem 1.4}
Lemma 1.4.1 If $A \subset B$ them $A\cap B = A$.\\ 
Lemma 1.4.2 If $A \subset B$ them $A\cup B = B$.\\ 
Lemma 1.4.3
$C\setminus B = C \cap \bar{B}$\\
\begin{enumerate}[label=\alph*)]
\item
$ (C\setminus B) \subset (C\setminus A ) \implies  (C\setminus B ) \cap  (C\setminus A )=  (C\setminus B )$\\
$(C\setminus B ) \cap  (C\setminus A ) = (C\cap \bar{B} ) \cap  (C \cap \bar{A} ) = C\cap \bar{B} \cap  C \cap \bar{A} = C\cap (\bar{B} \cap \bar{A})$ \\
that because union is asociative and comunative and we have the property $C\cap C = C$.\\
From DeMorgan property we have $\bar{B} \cap \bar {A}=\overline{B\cup A}$. \\From Lemma 1.4.2  $\overline{B\cup A} =  \overline{B} \Rightarrow C\cap (\bar{B} \cap \bar{A})= C \cap \bar{B}$.
 And using DeMorgan property we get $C \cap \bar{B} = C\setminus B$

 \item $(A\setminus B)\setminus(B\setminus C)= A \setminus B$ \\
 $(A\setminus B)\setminus(B\setminus C) = (A\cap \bar{B})\cap(\overline{B\cap \bar{C}})= (A\cap \bar{B})\cap({\bar{B}\cup \bar{\bar{C}}})$\\
 we have $\bar{\bar{C}} = C$ because the set contain all the elelemnts that are not contained in the set that has all the element not contained in set C is C.\\
 $(A\cap \bar{B})\cap(\bar{B}\cup C)= (A\cap \bar{B}\cap\bar{B})\cup (A\cap \bar{B}\cap C) = (A\cap \bar{B})\cup (A\cap \bar{B}\cap C) $ \\
 By $E$ we define the set containing evrything\\
$(A\cap \bar{B})\cup (A\cap \bar{B}\cap C) = ((A\cap \bar{B}) \cap E)\cup ((A\cap \bar{B})\cap C) = (A\cap \bar{B}) \cap (E \cup C)$\\
From the property on the evrything set $
\begin{cases}
  E \cup C = E \\
  E \cap C = C
\end{cases} \Rightarrow \\ \Rightarrow (A\cap \bar{B}) \cap (E \cup C) = (A \cap \bar{B})\cap E =A \cap \bar{B} = A \setminus B$

\item $A \cup (A \cap B) = A$\\
we allredy did this at point b) but i will do it once again here. 
By $E$ we define the set containing evrything\\
From the property on the evrything set $
\begin{cases}
  E \cup A = E \\
  E \cap A = A
\end{cases}$
\\
$A \cup (A \cap B) = (A \cap E) \cup (A \cap B) = A \cap (E \cup B) = A$

\end{enumerate}

\section*{Problem 1.5}
\[\left(\bigcap_{i=1}^n A_i\right)^c=\bigcup_{i=1}^n A_i^c\]\\
If $X$ is an element of all sets $A_i$ them $x \in  \bigcap_{i=1}^n A_i \Rightarrow x \not\in \left(\bigcap_{i=1}^n A_i\right)^c$ and olso that means that $X$ is in none of the sets $A_i^c$ sets them $X \not\in \bigcup_{i=1}^n A_i^c$ and both conditions are equivalent wich implies that DeMorgan indenty rule is true for $X$ in all sets 
$A_i$.\\
If $X$ is not an element of any set $A_i$ them $x \not\in  \bigcap_{i=1}^n A_i \Rightarrow x \in \left(\bigcap_{i=1}^n A_i\right)^c$ and olso that means that $X$ is in at least one of the sets $A_i^c$ them $X \not\in \bigcup_{i=1}^n A_i^c$ and both conditions are equivalent wich implies that DeMorgan indenty rule is true for $X$ in none of the sets 
$A_i$.

\section*{Problem 1.6}
\begin{enumerate}[label=\alph*)]
    \item $X \times Y = \{(a, b), (a, c), (a, d), (b, b), (b, c), (b, d)\}$
    \item $Y^2 = Y \times Y = \{(b, b), (b, c), (b, d), (c, b), (c, c), (c, d), (d, b), (d, c), (d, d)\}$
    \item $Y \times X \times X = \{(b, a, a), (b, a, b), (b, b, a), (b, b, b), (c, a, a), (c, a, b), (c, b, a), (c, b, b),$\\$ (d, a, a), (d, a, b), (d, b, a), (d, b, b)\}$
\end{enumerate}
At  point $c)$ we just treat $Y \times X$ as a new set.

\section*{Problem 1.7}
$A \times B = \{(a, b) \;|\; a \in A,\; b \in B\}$
If $A \times B = B \times A$ that means that any arbitrary orderet set $(a,\,b)$ should exist in both producs. So we have $(a,\,b) \in A \times B$ for $a \in A, \; b \in B$ and \\ $(a,\,b) \in B \times A$ for $b \in A, \; a \in B$ wich results that $a$ is in $A$ and $B$ and $b$ is in $A$ and $B$ and both $a$ and $b$ are any arbitraly element from $A$ and conseuncutivly $B$ wich means that this is posible only if $A=B$.
 If $A \times B = \varnothing$ either $A$ or $B$ should be $varnothing$ because from the property of the cartesian product if $a \in \varnothing$ there are no ordered sets to pair because $\varnothing$ is an empty set. Consequancely   $B \times A = \{(a, b) \;|\; a \in B,\; b \in A\}$ and is again $\varnothing$ because $A=\varnothing$ so there does not exist any $b$ element so we proved that 
 $
\begin{cases}
  A \times B = \varnothing \\
  B \times A = \varnothing
\end{cases} \rightarrow A \times B = B \times A = \varnothing
 $
\newpage
 \section*{Problem 1.8}
 \begin{enumerate}[label=\arabic*.]
    \item $f(n) = a$\\
    $n+2=a \Rightarrow n = 2-a$ this is a straing line so it only maps to one element so is injective
    $f(0)=0+2=2$,  $f(n)=0 \Rightarrow 2+n = 0 \Rightarrow n = -2 \not \in \ \mathbb{N}$ wich means is not surjective and bijective
    \item This is another straing line so is injective because it has only one solution for any $g(x) = a$. It has a interesting properti in only outputs odd numbers so is not surjective and bijective.\\
    $g(0) = -5$, $g(n) = 0 \Rightarrow 2n-5 = 0 \Rightarrow n = 2.5 \not\in \mathbb{Z}$
    \item we define $h(x)$ as 
    $h(x) = \sum_{k=0}^{\lfloor \log_{10} x \rfloor} \left( \lfloor x \cdot 10^{-k} \rfloor \bmod 10 \right)$ where $\lfloor x \rfloor$ is the lower bound of x. $h(21) = 3$ and $h(12) = 3$ so $h(x)$ is not injective. $h(x)$ can input negative numbers but do to this nature cann't output negative numbers because the digits of a number are hole numbers betwen 0 and 9 so $h(x)$ is not surjective or bijective
    \item Lets find the function minimum. $q'(n)=(n^2-n)'=2n-1$ so the minimum is at point $n=0.5$ so if we compute the adiacent point $q(0)=0$, $q(1)=0$ so is not injective. \\
    Second method $n^2-n=n(n-1)=0 \Rightarrow \begin{cases}  
  n &= 0 \\
  n&= 1
\end{cases}$ and because we have 2 solutions it is not injective.
It is neigher surjective or bijective because $q(2)=2$ and $q(3)=6$ but $q(n) = 3 \Rightarrow n=\frac{1\pm\sqrt{13}}{2}$ and $n\not\in \mathbb{Z}$
\item
    $r(n_1) = r(n_2) \Rightarrow \frac{1}{n_1} = \frac{1}{n_2} \Rightarrow n_1 = n_2$ so it maps distinct elements to distinct outputs, which means is injective.\\
    For $q = 2 \in \mathbb{Q}^+$, $r(n) = 2 \Rightarrow \frac{1}{n} = 2 \Rightarrow n = 0.5 \not\in \mathbb{N}^+$ (also $2/3 \in \mathbb{Q}^+$ has no integer pre-image) wich means is not surjective and bijective.

\end{enumerate}

\section*{Problem 1.9}

\begin{enumerate}[label=\alph*)]

    \item \textbf{Finite.}
    
    By the Fundamental Theorem of Algebra, a polynomial equation of degree
    $5$ has at most $5$ roots in $\mathbb{C}$. Therefore, the set of
    solutions is finite.

    \item \textbf{Countably infinite.}
    
    The set is
    \[
    S = \{n^2 \mid n \in \mathbb{N}\}.
    \]
    The function $f(n) = n^2$ defines a bijection
    $f:\mathbb{N}\to S$. Therefore, $S$ is countably infinite.

    \item \textbf{Uncountable.}
    
    $\mathbb{C}$ is in bijection with $\mathbb{R}^2$. Since $\mathbb{R}$
    is uncountable, $\mathbb{R}^2$ is also uncountable. Therefore,
    $\mathbb{C}$ is uncountable.

    \item \textbf{Finite.}
    
    Regardless of the actual number of Romanian words, every word of
    length at most $10$ can be formed using the finite alphabet of
    $31$ letters. Thus, the total number of possible strings of length
    at most $10$ is bounded by
    \[
    \sum_{k=1}^{10} 31^k < \infty.
    \]
    Therefore, the set of Romanian words with no more than $10$ letters
    is finite.

    \item \textbf{Depends on the sets $A$ and $B$.}
    
    \begin{itemize}
        \item If $|A| \neq |B|$, then there are no bijections from $A$
        to $B$, so the set of bijections is empty and therefore finite.

        \item If $A$ and $B$ are finite with $|A| = |B| = n$, then there
        are exactly $n!$ bijections, so the set is finite.

        \item If $A$ and $B$ are infinite with $|A| = |B|$, then the set
        of bijections is uncountable.
    \end{itemize}

    \item \textbf{Uncountable.}
    
    Since $B \subseteq A \cup B$ and $B$ is uncountable, the set
    $A \cup B$ must also be uncountable.

    \item \textbf{Uncountable.}
    
    Lines with a fixed slope $a$ are given by
    \[
    y = ax + b,
    \]
    where $b \in \mathbb{R}$. Each such line is uniquely determined by
    its intercept $b$, giving a bijection with $\mathbb{R}$. Therefore,
    the set of all lines with the same slope is uncountable.

    \item \textbf{Countably infinite.}
    
    Every English word can be represented as a finite string over a
    finite alphabet. The set of all finite strings over a finite alphabet
    is countably infinite. Therefore, the set of all English words is
    countably infinite.

    \item \textbf{Uncountable.}
    
    Every circle in $\mathbb{R}^2$ is uniquely determined by its center
    $(x_0,y_0) \in \mathbb{R}^2$ and its positive radius $r \in \mathbb{R}^+$.
    Since $\mathbb{R}^2 \times \mathbb{R}^+$ is uncountable, the set of
    all circles is uncountable.

    \item \textbf{Finite.}
    
    The set contains exactly
    \[
    2^{10^{10}}
    \]
    elements, which is a fixed finite integer. Therefore, the set is
    finite.

    \item \textbf{Countably infinite.}
    
    The set of all finite bit strings is
    \[
    \{0,1\}^* = \bigcup_{n=0}^{\infty} \{0,1\}^n.
    \]
    Each set $\{0,1\}^n$ is finite, and a countable union of finite sets
    is countable. Since there are infinitely many finite bit strings,
    the set is countably infinite.

\end{enumerate}


\section*{Problem 1.10}

\begin{enumerate}[label=\alph*)]

    \item \textbf{Finite.}
    
    Let
    \[
    A = [0,1], \qquad B = [1,2].
    \]
    Both $A$ and $B$ are uncountable intervals in $\mathbb{R}$.
    However,
    \[
    A \cap B = \{1\},
    \]
    which is finite.

    \item \textbf{Countably infinite.}
    
    Let
    \[
    A = (0,1) \cup \mathbb{Z}, \qquad
    B = (2,3) \cup \mathbb{Z}.
    \]
    Both $A$ and $B$ are uncountable. Their intersection is
    \[
    A \cap B = \mathbb{Z},
    \]
    which is countably infinite.

    \item \textbf{Uncountable.}
    
    Let
    \[
    A = [0,2], \qquad B = [1,3].
    \]
    Both $A$ and $B$ are uncountable. Their intersection is
    \[
    A \cap B = [1,2],
    \]
    which is an interval in $\mathbb{R}$ and therefore uncountable.

\end{enumerate}


\section*{Problem 1.11}

Assume by contradiction that $\mathbb{N}^{\infty}$ is countable. Then its
elements can be enumerated in a list
$S_1,S_2,S_3,\dots$, where each $S_i$ is an infinite sequence of
non-negative integers:

\begin{align*}
S_1 &= (s_{1,1},\,s_{1,2},\,s_{1,3},\,\dots) \\
S_2 &= (s_{2,1},\,s_{2,2},\,s_{2,3},\,\dots) \\
S_3 &= (s_{3,1},\,s_{3,2},\,s_{3,3},\,\dots) \\
&\;\vdots
\end{align*}

We construct a sequence
\[
D = (d_1,d_2,d_3,\dots)
\]
and define each term by
\[
d_k = s_{k,k} + 1
\qquad \text{for all } k \geq 1.
\]

Since $s_{k,k} \in \mathbb{N}$, it follows that $d_k \in \mathbb{N}$
for all $k \geq 1$. Thus,
\[
D \in \mathbb{N}^{\infty}.
\]

However, $D$ differs from $S_k$ at the $k$-th entry, since
\[
d_k = s_{k,k} + 1 \neq s_{k,k}.
\]
Therefore,
\[
D \neq S_k
\]
for every $k \geq 1$. Consequently, $D$ is not in the enumeration.

This contradicts the assumption that the list contains all elements of
$\mathbb{N}^{\infty}$. Therefore, $\mathbb{N}^{\infty}$ is uncountable.

\end{document}